\section{Множества и логические операции над ними}

Логические операции над множествами в настоящее время применяются не только
в математике, но и в многочисленных ее приложениях. В данной работе им
отводится особая роль: формализация, достигаемая на уровне алгебры
множеств, будет представлять собой существенное звено в конструктивных
построениях, осуществляемых с помощью $R$-функций (в частности, при
построении уравнений сложных геометрических объектов и структур решений
краевых задач). Заметим, что в ряде случаев наряду с термином
<<множество>> используются его синонимы <<совокупность>>,
<<многообразие>>, <<система>> и~др.

Примем следующие обозначения:

\begin{itemize}[nosep,label={},leftmargin=1em]
\item $\N$ "--- множество неотрицательных целых чисел;
\item $\Z$ "--- множество произвольных целых чисел;
\item $\Q$ "--- множество рациональных чисел;
\item $\R$ "--- множество действительных чисел;
\item $C^m(\Omega)$ "--- множество $m$ раз непрерывно дифференцируемых функций;
\item $\emptyset$ "--- пустое множество;
\item $2^A$ "--- множество всех подмножеств множества $A$;
\item $f(A)$ "--- образ множества $A$ при отображении $f$;
\item $f^{-1}(A)$ "--- прообраз множества $A$ при отображении $f$;\fixed{в оригинале <<при отображении $|$>>.}
\item $f/A$ "--- сужение (ограничение) отображения $f$ на множество $A$
  (т.\,е. $f/A$ есть отображение, совпадающее с $f$ на множестве $A$. Вне
  множества $A$ отображение $f/A$, вообще, не определено);
\item $\inf A$ ($\sup A$) "--- нижняя (верхняя) грань множества $A\subset\R$;
\item $a\in A$ (или $A\ni a$) "--- элемент $a$ принадлежит множеству $A$
  ($a$ находится в состоянии принадлежности к $A$);
\item $a\notin A$ (или $A\not\ni a$) "--- элемент $a$ не принадлежит
  множеству $A$ ($a$ находится в состоянии непринадлежности к $A$);
\item $A\subset B$ (или $B\supset A$) "--- множество $A$ является
  подмножеством множества $B$;
\item $A\equiv B$ "--- множества $A$ и $B$ тождественны, т.\,е.
  $A\supset B$ и $A\subset B$;
\item $\exists a\in A\dots$ "--- существует (по крайней мере один) элемент
  $a$ множества $A$, что \dots\ ($\exists$ "--- квантор существования);
\bookpage{14}%
\item $\forall a\in A\dots$ "--- для всякого элемента $a$ множества
  $A$\dots\ ($\forall$ "--- квантор общности);
\item $A_1\times\dots\times A_n$ "--- прямое произведение множеств
  $A_1,\dots,A_n$ есть множество упорядоченных строк $a_1,\dots,a_n$, где
  $a_i\in A_i$. В частности, если $A_1=\dots=A_n$, то
  $A\times\dots\times A=A^n$ называется $n$-й степенью множества $A$.
  Например, $n$-мерное евклидово пространство $\R^n$ есть $n$-я степень
  множества~$\R$.
\end{itemize}

Перейдем к логическим операциям над множествами. Операцию $F$ над
множествами $A,B,\dots$ будем называть логической, если в результате ее
применения получается множество, принадлежность элементов к которому
вполне определяется их состояниями принадлежности или непринадлежности
к множествам $A,B,\dots$ Примеры логических операций следующие:

1. $A\cap B$ (или $AB$) "--- \emph{пересечение} множеств $A$ и $B$ "--- есть
множество, состоящее из всех элементов, принадлежащих как множеству $A$,
так и множеству $B$.

2. $A\cup B$ (или $A+B$) "--- \emph{соединение} (сумма) множеств $A$ и $B$
"--- есть множество, состоящее из всех элементов, принадлежащих хотя бы
одному из множеств $A$ или $B$.

3. $A\setminus B$ (или $A-B$) "--- \emph{разность} множеств $A$ и $B$ "---
есть совокупность тех элементов $A$, которые не принадлежат $B$:
\begin{equation}\tag{1.1}
A\setminus B=A\cap\ov B.
\end{equation}

4. $A\sd B$ "--- \emph{симметрическая разность} "--- есть соединение
разностей $A\setminus B$ и $B\setminus A$, т.\,е.
\begin{equation}\tag{1.2}
A\sd B=(A\cap\ov B)\cup(B\cap\ov A).
\end{equation}

Часто приходится рассматривать системы множеств, являющиеся подмножествами
некоторого основного (так называемого унитарного) множества $I$. Для них
вводится понятие дополнения (или отрицания).

5. Множество $\ov A=I\setminus A$ называется дополнением (отрицанием)
множества $A$. (Иногда употребляются обозначения $\ov A=A'={\sim}A$ и~др.)

6. Тождество $F(A,B,\dots)\equiv\Phi(A,B,\dots)$, где $F$ и $\Phi$ "---
некоторые формулы, построенные с помощью тех или иных логических операций
над множествами $A,B,\dots$, означает равенство множеств $F(A,B,\dots)$
и $\Phi(A,B,\dots)$ при любом выборе множеств $A,B,\dots$

Нетрудно убедиться, что операции пересечения, соединения и дополнения
обладают следующими свойствами, выраженными в виде тождеств.

I. Коммутативность:
\[
1^\circ)\ A\cap B\equiv B\cap A;\qquad 2^\circ)\ A\cup B\equiv B\cup A.
\]
II. Ассоциативность:
\[
\begin{aligned}
&3^\circ)\ A\cap(B\cap C)\equiv(A\cap B)\cap C,\\
&4^\circ)\ A\cup(B\cup C)\equiv(A\cup B)\cup C.
\end{aligned}
\]
\bookpage{15}%
III. Дистрибутивность:
\[
\begin{aligned}
&5^\circ)\ A\cap(B\cup C)\equiv(A\cap B)\cup(A\cap C);\\
&6^\circ)\ A\cup(B\cap C)\equiv(A\cup B)\cap(A\cup C).
\end{aligned}
\]
IV. Идемпотентность:
\[
7^\circ)\ A\cap A\equiv A;\qquad 8^\circ)\ A\cup A\equiv A.
\]
V. Закон двойного отрицания:
\[
9^\circ)\ \ov{(\ov A)}\equiv\ov{\ov A}\equiv A.
\]
VI. Правила де Моргана:
\[
10^\circ)\ \ov{A\cap B}\equiv\ov A\cup\ov B;\qquad
11^\circ)\ \ov{A\cup B}\equiv\ov A\cap\ov B.
\]
VII. Закон логического противоречия:
\[
12^\circ)\ A\cap\ov A\equiv\emptyset.
\]
VIII. Закон исключенного третьего:
\[
13^\circ)\ A\cup\ov A\equiv I.
\]
IX. Операции с унитарным множеством $I$ и пустым множеством $\emptyset$:
\[
\begin{array}{ll}
14^\circ)\ I\cap A\equiv A; & 17^\circ)\ \emptyset\cup A\equiv A;\\
15^\circ)\ I\cup A\equiv I; & 18^\circ)\ \ov\emptyset\equiv I;\\
16^\circ)\ \emptyset\cap A\equiv\emptyset; & 19^\circ)\ \ov I\equiv\emptyset.
\end{array}
\]
Приведенные формулы позволяют осуществлять тождественные преобразования
формул, записанных в символах операций $A\cap B$, $A\cup B$, $\ov A$,
например, с целью их упрощения.

{\small\noindent\textbf{Пример.}
\[
(A\cap\ov B\cap C)\cup(A\cap C)\equiv(A\cap C\cap\ov B)\cup(A\cap C\cap I)
\equiv(A\cap C)\cap(\ov B\cup I)\equiv(A\cap C)\cap I\equiv A\cap C.
\]
(Здесь использованы свойства $1^\circ$, $3^\circ$, $5^\circ$, $14^\circ$,
$15^\circ$.)\fixed{в оригинале в левой части стоит
$(A\cap\ov B\cap C)\cap(A\cap C)$; по ходу преобразований (вынесение
$A\cap C$ по закону $5^\circ$) первым знаком должен быть $\cup$.}\par}

Логическая терминология, использованная для наименования некоторых из
приведенных выше формул, имеет свои корни в метаматематике~[84],
в частности в исчислении высказываний или булевой алгебре. О последней
речь будет идти подробнее в \S\S\,6--8 данной главы.

Справедливость свойств $1^\circ$--$19^\circ$ легко проверяется. Убедимся,
например, в справедливости формулы $5^\circ$. Пусть
$x\in A\cap(B\cup C)$, т.\,е. $x$ "--- элемент множества, стоящего в левой
части формулы $5^\circ$. Это означает, что $x$ принадлежит одновременно как
множеству $A$, так и множеству $B\cup C$. Но тогда $x\in B$ или $x\in C$.
Следовательно, $x$ принадлежит по крайней мере одному из множеств
$A\cap B$ или $A\cap C$. Тогда $x\in(A\cap B)\cup(A\cap C)$ и, таким
образом, $A\cap(B\cup C)\subset(A\cap B)\cup(A\cap C)$. Наоборот, если
$x\in(A\cap B)\cup(A\cap C)$, то $x\in A\cap B$ или $x\in A\cap C$.
Следо\-%
\bookpage{16}%
вательно, $x\in A$ и по крайней мере одному из множеств $B$ или $C$. Но
тогда $x\in B\cup C$, а значит, и $x\in A\cap(B\cup C)$. Таким образом,
$(A\cap B)\cup(A\cap C)\subset A\cap(B\cup C)$. Сравнивая эту формулу
с полученной ранее формулой $A\cap(B\cup C)\subset(A\cap B)\cup(A\cap C)$,
приходим к формуле~$5^\circ$.

Благодаря тому что операции $A\cap B$ и $A\cup B$ ассоциативны, при
многократном их применении можно опускать скобки, например
$((A\cap B)\cap C)\cap D\equiv A\cap B\cap C\cap D$. Применяется также
сокращенная запись:
\begin{equation}\tag{1.3}
A_1\cap\dots\cap A_n\equiv\bigcap_{i=1}^{i=n}A_i;\qquad
A_1\cup\dots\cup A_n\equiv\bigcup_{i=1}^{i=n}A_i.
\end{equation}

\subsection*{Упражнения и задачи}
\label{ex:s1}
1. Доказать тождества
\begin{equation}\tag{1.4}
\begin{gathered}
A\sd B\equiv(A\cup B)\setminus(A\cap B);\\
A\cup B\equiv(A\sd B)\sd(A\cap B);\qquad A\setminus B\equiv A\sd(A\cap B);\\
(A\cap B)\cup B\equiv B;\qquad (A\cup B)\cap B\equiv B;\\
(A\cap B)\cup(A\cap\ov B)\equiv A;\qquad (A\cup B)\cap(A\cup\ov B)\equiv A;\\
(\ov A\cap B)\cup(A\cap\ov B)\equiv(A\cup B)\cap(\ov A\cup\ov B).
\end{gathered}
\end{equation}
(Для доказательства использовать формулы (1.1), (1.2) и свойства
$1^\circ$--$19^\circ$.)

2. Показать, что симметрическая разность "--- это коммутативная
и ассоциативная операции.

\textsc{Решение.} В соответствии с формулой (1.2) и свойствами
$1^\circ$--$19^\circ$ получаем
\begin{align*}
A\sd(B\sd C)&\equiv[A\cap(\ov{B\sd C})]\cup[\ov A\cap(B\sd C)]\equiv\\
&\equiv\{A\cap\ov{[(B\cap\ov C)\cup(\ov B\cap C)]}\}\cup\{\ov A\cap[(B\cap\ov C)\cup(\ov B\cap C)]\}\equiv\\
&\equiv\{A\cap[\ov{(B\cap\ov C)}\cap\ov{(\ov B\cap C)}]\}\cup\{\ov A\cap[(B\cap\ov C)\cup(\ov B\cap C)]\}\equiv\\
&\equiv\{A\cap[(\ov B\cup C)\cap(B\cup\ov C)]\}\cup\{(\ov A\cap B\cap\ov C)\cup(\ov A\cap\ov B\cap C)\}\equiv\\
&\equiv\{[A\cap(\ov B\cup C)]\cap[A\cap(B\cup\ov C)]\}\cup(\ov A\cap B\cap\ov C)\cup{}\\
&\qquad\cup(\ov A\cap\ov B\cap C)\equiv\{[(A\cap\ov B)\cup(A\cap C)]\cap[(A\cap B)\cup(A\cap\ov C)]\}\cup{}\\
&\qquad\cup(\ov A\cap B\cap\ov C)\cup(\ov A\cap\ov B\cap C)\equiv\\
&\equiv\{[(A\cap\ov B)\cap(A\cap B)]\cup[(A\cap\ov B)\cap(A\cap\ov C)]\cup[(A\cap C)\cap(A\cap B)]\cup{}\\
&\qquad\cup[(A\cap C)\cap(A\cap\ov C)]\}\cup(\ov A\cap B\cap\ov C)\cup(\ov A\cap\ov B\cap C)\equiv\\
&\equiv\{[(A\cap A)\cap(B\cap\ov B)]\cup[(A\cap A)\cap(\ov B\cap\ov C)]\cup[(A\cap A)\cap(C\cap B)]\cup{}\\
&\qquad\cup[(A\cap A)\cap(C\cap\ov C)]\}\cup(\ov A\cap B\cap\ov C)\cup(\ov A\cap\ov B\cap C)\equiv\\
&\equiv(A\cap\emptyset)\cup(A\cap\ov B\cap\ov C)\cup(A\cap B\cap C)\cup(A\cap\emptyset)\cup{}\\
&\qquad\cup(\ov A\cap B\cap\ov C)\cup(\ov A\cap\ov B\cap C)\equiv\emptyset\cup(A\cap\ov B\cap\ov C)\cup{}\\
&\qquad\cup(A\cap B\cap C)\cup\emptyset\cup(\ov A\cap B\cap\ov C)\cup(\ov A\cap\ov B\cap C)\equiv\\
&\equiv(A\cap\ov B\cap\ov C)\cup(A\cap B\cap C)\cup(\ov A\cap B\cap\ov C)\cup(\ov A\cap\ov B\cap C).
\end{align*}
\bookpage{17}Поскольку последнее выражение инвариантно относительно циклической
перестановки букв (т.\,е. при замене $A$, $B$, $C$ на $C$, $A$, $B$
соответственно), то оно равно также $C\sd(A\sd B)$. Таким образом,
ассоциативность операции $A\sd B$ доказана. Ее коммутативность
очевидна.\fixed{в оригинале в восьмой строке выкладки стоит скобка
<<$([A\cap A)$>> вместо <<$[(A\cap A)$>>.}

3. Пусть $I=\R$. Показать, что равенство
$(A\cap B)\cup C=C\cup(B\cap C)$ справедливо, если $A=[a,b]$, $B=[c,d]$,
$C=[e,f]$ и $e\le\max(a,c)$; $f\ge\min(b,d)$, но оно не является
тождеством.

4. Показать, что правила де Моргана $10^\circ$ и $11^\circ$ верны для
$n$ ($n\ge2$) компонент, т.\,е.
\begin{equation}\tag{1.5}
\ov{\bigcap_{i=1}^{i=n}A_i}\equiv\bigcup_{i=1}^{i=n}\ov{A_i};\qquad
\ov{\bigcup_{i=1}^{i=n}A_i}\equiv\bigcap_{i=1}^{i=n}\ov{A_i}.
\end{equation}

%======================================================================
\begin{solution}[1 (\S\,1)]
Будем пользоваться только (1.1), (1.2) и свойствами $1^\circ$--$19^\circ$.
Полезно заранее заметить два следствия.

(а) \emph{Поглощение:} $(A\cap B)\cup B\equiv B$ и $(A\cup B)\cap B\equiv B$.
Действительно,
$(A\cap B)\cup B\overset{14^\circ}{\equiv}(A\cap B)\cup(I\cap B)
\overset{5^\circ}{\equiv}(A\cup I)\cap B\overset{15^\circ}{\equiv}I\cap B\equiv B$;
второе тождество двойственно первому (те же шаги с $6^\circ$, $17^\circ$,
$16^\circ$):
$(A\cup B)\cap B\equiv(A\cup B)\cap(\emptyset\cup B)\equiv(A\cap\emptyset)\cup B\equiv\emptyset\cup B\equiv B$.
Это четвертое и пятое тождества (1.4).

(б) \emph{Склеивание:}
$(A\cap B)\cup(A\cap\ov B)\overset{5^\circ}{\equiv}A\cap(B\cup\ov B)
\overset{13^\circ}{\equiv}A\cap I\equiv A$;
$(A\cup B)\cap(A\cup\ov B)\overset{6^\circ}{\equiv}A\cup(B\cap\ov B)
\overset{12^\circ}{\equiv}A\cup\emptyset\equiv A$. Это шестое и седьмое тождества.

\smallskip
\emph{Последнее тождество.} Раскроем скобки в правой части по $5^\circ$:
\[
(A\cup B)\cap(\ov A\cup\ov B)\equiv(A\cap\ov A)\cup(A\cap\ov B)\cup(B\cap\ov A)\cup(B\cap\ov B)
\equiv(A\cap\ov B)\cup(\ov A\cap B),
\]
так как $A\cap\ov A\equiv B\cap\ov B\equiv\emptyset$. Попутно получена
полезная формула $A\sd B\equiv(A\cup B)\cap(\ov A\cup\ov B)$.

\emph{Первое тождество.} По (1.1) и $10^\circ$:
$(A\cup B)\setminus(A\cap B)\equiv(A\cup B)\cap\ov{A\cap B}\equiv(A\cup B)\cap(\ov A\cup\ov B)$,
а это, как только что показано, $A\sd B$.

\emph{Третье тождество.} По (1.2), $10^\circ$ и $9^\circ$:
\begin{multline*}
A\sd(A\cap B)\equiv\bigl(A\cap\ov{A\cap B}\bigr)\cup\bigl(\ov A\cap A\cap B\bigr)
\equiv\bigl(A\cap(\ov A\cup\ov B)\bigr)\cup\emptyset\equiv{}\\
\equiv(A\cap\ov A)\cup(A\cap\ov B)\equiv A\cap\ov B\equiv A\setminus B.
\end{multline*}

\emph{Второе тождество.} Обозначим $D=A\sd B=(A\cap\ov B)\cup(\ov A\cap B)$,
$P=A\cap B$. Заметим, что $D\cap P\equiv\emptyset$ (каждый член $D$
содержит $\ov A$ или $\ov B$). Тогда $D\cap\ov P\equiv D$ (так как
$D\equiv(D\cap P)\cup(D\cap\ov P)$ по склеиванию) и аналогично
$P\cap\ov D\equiv P$. Поэтому
\[
(A\sd B)\sd(A\cap B)\equiv(D\cap\ov P)\cup(\ov D\cap P)\equiv D\cup P
\equiv(A\cap\ov B)\cup(\ov A\cap B)\cup(A\cap B).
\]
Склеивая $(A\cap\ov B)\cup(A\cap B)\equiv A$ и $(\ov A\cap B)\cup(A\cap B)\equiv B$
(член $A\cap B$ можно повторить по $8^\circ$), получаем $A\cup B$.
\end{solution}

\begin{solution}[2 (\S\,1): дополнение]
Решение в книге приводит $A\sd(B\sd C)$ к виду
$K_{100}\cup K_{111}\cup K_{010}\cup K_{001}$, где
$K_{\alpha\beta\gamma}=A^\alpha\cap B^\beta\cap C^\gamma$ ($X^1=X$,
$X^0=\ov X$). Это ровно те <<кусочки>> диаграммы Эйлера, в которых точка
принадлежит \emph{нечетному} числу множеств. Отсюда сразу видна и
коммутативность, и ассоциативность: для любых $A_1,\dots,A_n$ множество
$A_1\sd\dots\sd A_n$ (при любой расстановке скобок) "--- это множество
точек, принадлежащих нечетному числу из $A_i$. (Индукция: $x\in X\sd Y$
$\iff$ $x$ лежит ровно в одном из $X$, $Y$, т.\,е. число вхождений в $X\sd Y$
имеет ту же четность, что сумма чисел вхождений.) Коммутативность $A\sd B\equiv B\sd A$
следует из (1.2) и $1^\circ$, $2^\circ$.
\end{solution}

\begin{solution}[3 (\S\,1)]
По поглощению $C\cup(B\cap C)\equiv C$, поэтому равенство
$(A\cap B)\cup C=C\cup(B\cap C)$ равносильно $(A\cap B)\cup C=C$, т.\,е.
\[
A\cap B\subset C .
\]
Для отрезков $A\cap B=[\max(a,c),\ \min(b,d)]$ (пусто, если
$\max(a,c)>\min(b,d)$). Если $e\le\max(a,c)$ и $f\ge\min(b,d)$, то
$[\max(a,c),\min(b,d)]\subset[e,f]$, и равенство выполнено (если же
$A\cap B=\emptyset$, оно тем более верно).

Тождеством равенство не является: при $A=B=[0,1]$, $C=[2,3]$ слева
$[0,1]\cup[2,3]$, справа $[2,3]$. Вообще, равенство верно в точности при
$A\cap B\subset C$; условия задачи "--- достаточные (а при
$A\cap B\neq\emptyset$ и необходимые).
\end{solution}

\begin{solution}[4 (\S\,1)]
Индукция по $n$. При $n=2$ это $10^\circ$, $11^\circ$. Пусть (1.5) верно
для $n$. Тогда по ассоциативности, $10^\circ$ и предположению индукции
\[
\ov{\bigcap_{i=1}^{n+1}A_i}\equiv\ov{\Bigl(\bigcap_{i=1}^{n}A_i\Bigr)\cap A_{n+1}}
\equiv\ov{\bigcap_{i=1}^{n}A_i}\cup\ov{A_{n+1}}
\equiv\Bigl(\bigcup_{i=1}^{n}\ov{A_i}\Bigr)\cup\ov{A_{n+1}}\equiv\bigcup_{i=1}^{n+1}\ov{A_i}.
\]
Второе тождество получается так же из $11^\circ$ (или из первого,
примененного к $\ov{A_i}$, с последующим взятием дополнения и законом
$9^\circ$). Можно рассуждать и непосредственно: $x\notin\bigcap A_i$
$\iff$ $x\notin A_i$ хотя бы для одного $i$ $\iff$ $x\in\bigcup\ov{A_i}$.
\end{solution}
